\chapter{คู่มือการติดตั้งและสอบเทียบระบบเซนเซอร์วัดดินและสภาพอากาศ}
\label{app:appendixA}

\section{ข้อกำหนดมาตรฐานการติดตั้งโพรบวัดความชื้นดิน}
1. ควรติดตั้งเซนเซอร์ให้ห่างจากหัวจ่ายน้ำสปริงเกลอร์ในระยะประมาณ 1.0 ถึง 1.5 เมตร เพื่อวัดปริมาณน้ำเฉลี่ยของทรงพุ่ม
2. ต้องหลีกเลี่ยงการติดตั้งใกล้รอยแตกของดินหรือบริเวณที่มีหินขนาดใหญ่ขวางกั้น
3. สายสัญญาณต้องร้อยผ่านท่อร้อยสายร้อยฝังดินลึกไม่น้อยกว่า 20 เซนติเมตร ป้องกันการถูกจอบหรือรถแทรกเตอร์ตัดขาด

\section{ขั้นตอนการสอบเทียบความถูกต้อง (Calibration Protocol)}
1. สอบเทียบจุดแห้งสนิท (Dry Point Calibration): อบตัวอย่างดินในตู้อบที่ 105$^{\circ}$C เป็นเวลา 24 ชั่วโมงจนแห้งสนิท เสียบโพรบวัดค่า VWC ต้องแสดงผล 0.0\%
2. สอบเทียบจุดอิ่มตัวด้วยน้ำ (Saturation Point Calibration): เติมน้ำจนท่วมดินอิ่มตัว วัดค่า VWC และเปรียบเทียบกับค่าความจุสนามจริง
